\section{Grenzwerte von Funktionen}
\label{sec:real_functions_limits}
Bisher haben wir Funktionen an vorgegebenen Stellen ausgewertet.
Jetzt ändern wir die Frage: Wie verhalten sich die Funktionswerte,
wenn sich die Eingabe einer Stelle immer weiter nähert?
% Handschrift: PDF-Seiten 19–20, Fall a) als Motivation.
\begin{ExampleBox}{Der Wert bei null und die Werte nahe null}{real_functions_limit_special_points}
  Wir betrachten die folgende stückweise definierte Funktion:
  \par\smallskip\noindent
  \begin{minipage}[c]{0.48\linewidth}
    \[
      f(x)=\begin{cases}1,&x\ne0,\\2,&x=0.\end{cases}
    \]
  \end{minipage}\hfill
  \begin{minipage}[c]{0.47\linewidth}
    \centering
    \CachedTikzFigure{content/resources/chapters/real_functions/tikz/functions_limit_special_points.tex}
  \end{minipage}
  \par\smallskip
  \begin{itemize}
    \item Der Funktionwert an der Stelle $x=2$ ist gegeben durch $f(2)=\Cloze{1cm}{1}.$
    \item Bei der Annäherung and $x \to 2$ von links und rechts (ohne $x=2$ zu erreichen)
    bleiben die Werte bei $1$. Wir sprechen vom Grenzwert und schreiben:
      \[
        \lim_{x\to2}f(x)=\Cloze{1cm}{1}.
      \]
  \end{itemize}
  \par\medskip
  Der Grenzwert (Annäherungswert) kann sich also vom Funktionswert unterscheiden.
\end{ExampleBox}

Um diese Unterscheidung zu formalisieren, führen wir den Begriff des
\emph{Grenzwerts} \EnglishTerm{limit} ein.
\begin{DefinitionBox}{Grenzwert -- anschauliche Bedeutung}{real_functions_limit}
  \noindent
  \begin{minipage}[c]{0.49\linewidth}
    Der \emph{Grenzwert} $L\in\R$ ist der Wert, dem sich die
    Funktionswerte $f(x)$ annähern, wenn $x\in D_f$ immer näher an $a$
    heranrückt, ohne $a$ selbst anzunehmen ($x\ne a$). Wir schreiben
    \[
      \lim_{x\to a}f(x)=L.
    \]
    Für einen definierten Grenzwert muss eine Annäherung von beiden Seiten zum
    selben Grenzwert $L$ führen.
  \end{minipage}\hfill
  \begin{minipage}[c]{0.47\linewidth}
    \centering
    \CachedTikzFigure{content/resources/chapters/real_functions/tikz/functions_limit_approach.tex}
  \end{minipage}
\end{DefinitionBox}

Wir verwenden hier einen anschaulichen Grenzwertbegriff, anstatt einer formalen Definition
über ein $\varepsilon$-$\delta$-Kriterium oder mittels konvergeter Folgen.

\begin{ExampleBox}{Grenzwert bei einer Sprungfunktion}{real_functions_jump_limit}
  Wir vergleichen die Annäherung an null bei der folgenden idealisierten
  Schaltfunktion:
  \par\smallskip
  \begin{minipage}[c]{0.47\linewidth}
    \[
      s(x)=\begin{cases}0,&x<0,\\1,&x\ge0.\end{cases}
    \]
    Die Annäherung an null liefert unterschiedliche Werte:
    \begin{itemize}
      \item von links ($x<0$): $\Cloze{1cm}{0}$,
      \item von rechts ($x>0$): $\Cloze{1cm}{1}$.
    \end{itemize}
  \end{minipage}\hfill
  \begin{minipage}[c]{0.49\linewidth}
    \centering
    \CachedTikzFigure{content/resources/chapters/real_functions/tikz/functions_jump.tex}
  \end{minipage}
  \par\medskip
  Die unterschiedlichen Annäherungswerte lassen keinen gemeinsamen
  Grenzwert zu:
  \[
    \lim_{x\to0}s(x)\qquad
    \Cloze{3.8cm}{\textcolor{red}{\text{existiert nicht.}}}
  \]
  Auch eine Änderung von $s(0)$ würde daran nichts ändern.
\end{ExampleBox}


\begin{ExampleBox}{Grenzwert trotz Definitionslücke}{real_functions_limit_domain_gap}
  Wir betrachten die Funktion $g(x)=x$ mit der Definitionsmenge
  $D_g=\R\setminus\{0\}$.
  \par\smallskip\noindent
  \begin{minipage}[c]{0.48\linewidth}
    \[
      \begin{aligned}
        g\colon D_g&\longrightarrow\R,\\
        x&\longmapsto x.
      \end{aligned}
    \]
    Der Grenzwert für $x\to0$ ist definiert, da sich die $y$-Werte
    von links und rechts derselben Zahl annähern:
    \[
      \lim_{x\to0}g(x)=\Cloze{1cm}{0}.
    \]
  \end{minipage}\hfill
  \begin{minipage}[c]{0.47\linewidth}
    \centering
    \CachedTikzFigure{content/resources/chapters/real_functions/tikz/functions_limit_domain_gap.tex}
  \end{minipage}
  \par\medskip
  Der Funktionswert $g(0)$ ist nicht definiert, nur der Grenzwert.
\end{ExampleBox}
%
\begin{KeyPointBox}[before skip=0.3\baselineskip]{Grenzwert trotz Definitionslücke}{real_functions_two_sided_criterion}
  Ist die Annäherung von beiden Seiten möglich, ist ein Grenzwert
  nur definiert, wenn beide Annäherungen denselben Wert liefern.
  Andernfalls ist der Grenzwert nicht definiert.
\end{KeyPointBox}

\begin{ExampleBox}{Eine endliche Stelle -- unbegrenzte Funktionswerte}{real_functions_reciprocal_zero}
  Wir vergleichen zwei Funktionen mit der Definitionsmenge
  $D_f=\R\setminus\{0\}$ bei der Annäherung an null.
  \[
    \lim_{x\to0}f(x).
  \]
  \par\smallskip\noindent
  \begin{minipage}[t]{0.47\linewidth}
    \begin{enumerate}[label=\alph*),leftmargin=*,topsep=0pt]
      \item \strut
    \end{enumerate}
    \vspace{-0.5\baselineskip}
    \begin{minipage}[c][5.2cm][c]{\linewidth}
      \centering
      \CachedTikzFigure{content/resources/chapters/real_functions/tikz/functions_reciprocal_square_limit.tex}
    \end{minipage}
    \par\smallskip
    Von beiden Seiten wachsen die Funktionswerte über jede Grenze:
    \[
      \lim_{x\to0}\frac1{x^2}=\Cloze{1.5cm}{+\infty}.
    \]
    Da $+\infty$ keine reelle Zahl ist, liegt kein endlicher Grenzwert vor.
    Man spricht von einem \emph{uneigentlichen Grenzwert}
    \EnglishTerm{infinite limit}.
  \end{minipage}\hfill
  \begin{minipage}[t]{0.47\linewidth}
    \begin{enumerate}[label=\alph*),start=2,leftmargin=*,topsep=0pt]
      \item \strut
    \end{enumerate}
    \vspace{-0.5\baselineskip}
    \begin{minipage}[c][5.2cm][c]{\linewidth}
      \centering
      \CachedTikzFigure{content/resources/chapters/real_functions/tikz/functions_reciprocal_infinity.tex}
    \end{minipage}
    \par\smallskip
    Die Funktionswerte gehen in entgegengesetzte Richtungen:
    \begin{itemize}[itemsep=0.2em]
      \item Von links: \Cloze{3.6cm}{unbegrenzt negativ}.
      \item Von rechts: \Cloze{3.6cm}{unbegrenzt positiv}.
    \end{itemize}
    \par\medskip
    Es gibt weder einen endlichen noch einen gemeinsamen
    uneigentlichen Grenzwert:
    \[
      \lim_{x\to0}\frac1x\quad
      \Cloze{3.4cm}{\textcolor{red}{\text{existiert nicht.}}}
    \]
  \end{minipage}
\end{ExampleBox}

\par\noindent\begin{minipage}{\linewidth}
\subsection{Grenzwerte im Unendlichen}
\label{subsec:real_functions_infinite_limits}
Häufig interessiert uns das asymptotische Verhalten einer Funktion,
also wie sie sich für betragsmäßig große $x$ verhält.
Dieses Verhalten beschreiben wir mit Grenzwerten für $x\to\pm\infty$.
\begin{DefinitionBox}{Grenzverhalten im Unendlichen}{real_functions_limit_at_infinity}
  Wächst $x$ unbegrenzt ins Positive (oder fällt unbegrenzt ins Negative),
  also $x\to+\infty$ (bzw. $x\to-\infty$), und nähern sich dabei die
  Funktionswerte $f(x)$ einem Wert $L\in\R$, so heißt $L$ der
  \emph{Grenzwert von $f$ im Unendlichen}. Wir schreiben:
  \[
    \lim_{x\to+\infty}f(x)=L
    \qquad\left(\text{bzw. }\lim_{x\to-\infty}f(x)=L\right).
  \]
\end{DefinitionBox}
\begin{ExampleBox}[before skip=0.3\baselineskip]{Grenzverhalten einer Geraden}{real_functions_identity_infinity}
  Am Graphen der Geraden betrachten wir das Verhalten für $x\to\pm\infty$.
  \[
    \lim_{x\to\pm\infty}f(x).
  \]
  \par\smallskip\noindent
  \begin{minipage}[c]{0.47\linewidth}
    Für $x\to+\infty$ wachsen die Funktionswerte unbegrenzt ins Positive;
    für $x\to-\infty$ fallen sie unbegrenzt ins Negative.
    \[
      \begin{aligned}
        \lim_{x\to+\infty}x&=\Cloze{1.6cm}{+\infty},\\[0.5em]
        \lim_{x\to-\infty}x&=\Cloze{1.6cm}{-\infty}.
      \end{aligned}
    \]
    Beide Grenzwerte sind uneigentlich.
  \end{minipage}\hfill
  \begin{minipage}[c]{0.47\linewidth}
    \centering
    \CachedTikzFigure{content/resources/chapters/real_functions/tikz/functions_identity_infinity.tex}
  \end{minipage}
\end{ExampleBox}
\end{minipage}

\subsection{Grenzwertgesetze}
\label{subsec:real_functions_limit_laws}
Grenzwerte müssen nicht jedes Mal durch eine neue Annäherungsbetrachtung
bestimmt werden. Wir zerlegen einen Term in bekannte Bausteine und
verknüpfen deren Grenzwerte.
\begin{TheoremBox}{Grenzwertgesetze}{real_functions_limit_laws}
  Existieren die Grenzwerte $\lim_{x\to a}f(x)=A$ und
  $\lim_{x\to a}g(x)=B$ mit $A,B\in\R$, so gelten:
  \[
    \begin{array}{l@{\qquad}l}
      \text{Summe/Differenz:}&
        \displaystyle\vphantom{\lim_{x\to a}\frac{f(x)}{g(x)}=\frac AB}
        \lim_{x\to a}\bigl(f(x)\pm g(x)\bigr)=A\pm B,\\[1em]
      \text{Produkt:}&
        \displaystyle\vphantom{\lim_{x\to a}\frac{f(x)}{g(x)}=\frac AB}
        \lim_{x\to a}f(x)g(x)=AB,\\[1em]
      \text{Quotient:}&
        \displaystyle\lim_{x\to a}\frac{f(x)}{g(x)}=\frac AB
        \qquad\text{falls }B\ne0.
    \end{array}
  \]
\end{TheoremBox}
\begin{ExampleBox}{Grenzverhalten von Potenzen}{real_functions_monomial_limit_laws}
  Mit den Grenzwertgesetzen bestimmen wir beispielhaft das asymptotische
  Verhalten der folgenden vier Funktionen für $x\to\pm\infty$.
  \par\smallskip\noindent
  \begin{minipage}[t]{0.47\linewidth}
    \begin{enumerate}[label=\alph*),leftmargin=*,topsep=0pt]
      \item Quadratfunktion
    \end{enumerate}
    \vspace{0.4\baselineskip}
    \begin{minipage}[c][5.2cm][c]{\linewidth}
      \centering
      \CachedTikzFigure{content/resources/chapters/real_functions/tikz/functions_square_infinity.tex}
    \end{minipage}
    \par\vspace{0.6\baselineskip}
    \begin{ClozeArea}[16]{7}
      \[
        \begin{aligned}
          \lim_{x\to\pm\infty}x^2&=\lim_{x\to\pm\infty}(x\cdot x)\\[1em]
          &=\underbrace{\lim_{x\to\pm\infty}x}_{\pm\infty}
            \cdot\underbrace{\lim_{x\to\pm\infty}x}_{\pm\infty}
          =\Cloze{1.6cm}{+\infty}.
        \end{aligned}
      \]
    \end{ClozeArea}
  \end{minipage}\hfill
  \begin{minipage}[t]{0.47\linewidth}
    \begin{enumerate}[label=\alph*),start=2,leftmargin=*,topsep=0pt]
      \item Kubikfunktion
    \end{enumerate}
    \vspace{0.4\baselineskip}
    \begin{minipage}[c][5.2cm][c]{\linewidth}
      \centering
      \CachedTikzFigure{content/resources/chapters/real_functions/tikz/functions_cubic_infinity.tex}
    \end{minipage}
    \par\vspace{0.6\baselineskip}
    \begin{ClozeArea}[16]{7}
      \[
        \begin{aligned}
          \lim_{x\to\pm\infty}x^3&=\lim_{x\to\pm\infty}(x^2\cdot x)\\[1em]
          &=\underbrace{\lim_{x\to\pm\infty}x^2}_{+\infty}
            \cdot\underbrace{\lim_{x\to\pm\infty}x}_{\pm\infty}
          =\Cloze{1.6cm}{\pm\infty}.
        \end{aligned}
      \]
    \end{ClozeArea}
  \end{minipage}
  \par\vspace{1.5\baselineskip}\noindent
  \begin{minipage}[t]{0.47\linewidth}
    \begin{enumerate}[label=\alph*),start=3,leftmargin=*,topsep=0pt]
      \item Hyperbel
    \end{enumerate}
    \vspace{0.4\baselineskip}
    \begin{minipage}[c][5.2cm][c]{\linewidth}
      \centering
      \CachedTikzFigure{content/resources/chapters/real_functions/tikz/functions_reciprocal_infinity.tex}
    \end{minipage}
    \par\vspace{0.6\baselineskip}
    \begin{ClozeArea}[16]{7}
    \[
      \lim_{x\to\pm\infty}\frac1x
      =\frac{\overbrace{\displaystyle\lim_{x\to\pm\infty}1}^{1}}
        {\underbrace{\displaystyle\lim_{x\to\pm\infty}x}_{\pm\infty}}
      =\Cloze{1cm}{0}.
    \]
    \end{ClozeArea}
  \end{minipage}\hfill
  \begin{minipage}[t]{0.47\linewidth}
    \begin{enumerate}[label=\alph*),start=4,leftmargin=*,topsep=0pt]
      \item Kehrwert der Quadratfunktion
    \end{enumerate}
    \vspace{0.4\baselineskip}
    \begin{minipage}[c][5.2cm][c]{\linewidth}
      \centering
      \CachedTikzFigure{content/resources/chapters/real_functions/tikz/functions_inverse_square_infinity.tex}
    \end{minipage}
    \par\vspace{0.6\baselineskip}
    \begin{ClozeArea}[16]{7}
    \[
      \begin{aligned}
        \lim_{x\to\pm\infty}\frac1{x^2}
        &=\lim_{x\to\pm\infty}\left(\frac1x\cdot\frac1x\right)\\[1em]
        &=\underbrace{\lim_{x\to\pm\infty}\frac1x}_{0}
          \cdot\underbrace{\lim_{x\to\pm\infty}\frac1x}_{0}
          =\Cloze{1cm}{0}.
      \end{aligned}
    \]
    \end{ClozeArea}
  \end{minipage}
\end{ExampleBox}
\Needspace{15\baselineskip}
Diese Beobachtungen verallgemeinern wir nun.
\begin{KeyPointBox}{Potenzen und ihre Kehrwerte im Unendlichen}{real_functions_power_limits}
  Für Potenzen $x^n$ mit ganzzahligem Exponenten $n$ gilt:
  \[
    \lim_{x\to\pm\infty}x^n=
    \begin{cases}
      \Cloze{1.2cm}{+\infty},&n>0,\ n\text{ gerade},\\[0.5em]
      \Cloze{1.2cm}{\pm\infty},&n>0,\ n\text{ ungerade},\\[0.5em]
      \Cloze{1.2cm}{1},&n=0,\\[0.5em]
      \Cloze{1.2cm}{0},&n<0.
    \end{cases}
  \]
\end{KeyPointBox}
Nachdem wir die Grenzwerte von Monomen bestimmt haben, können wir mit den
Grenzwertgesetzen auch die Grenzwerte von Polynomen bestimmen.
\begin{ExampleBox}{Höchste Potenz ausklammern}{real_functions_factored_polynomial_limit}
  Bestimmen Sie den Grenzwert von $p(x)=-3x^3+x^2$ für $x\to+\infty$.
  Zunächst versuchen wir, die Grenzwerte der einzelnen Summanden
  direkt zu verknüpfen:
  \begin{ClozeArea}[34]{3}
    \[
      \begin{aligned}
        \lim_{x\to+\infty}p(x)
          &\overset{?}{=} \lim_{x\to+\infty}(-3x^3) + \lim_{x\to+\infty}x^2
          \overset{?}{=}-\infty+\infty\quad
            \text{\raisebox{-.25\height}{\CachedTikzFigure{../_shared/resources/tikz/smiley_sad.tex}}}
      \end{aligned}
    \]
  \end{ClozeArea}
  Dabei handelt es sich um eine sogenannte \emph{unbestimmte Form},
  aus der sich der Grenzwert nicht
  bestimmen lässt.

  Stattdessen klammern wir zunächst die höchste Potenz
  aus und wenden dann Produkt- und Summenregel der Grenzwertgesetze an.
  \begin{ClozeArea}[34]{8}
    \[
      \begin{aligned}
        \lim_{x\to+\infty}p(x)
          &=\lim_{x\to+\infty}\left(-3x^3+x^2\right)\\[1em]
          &=\lim_{x\to+\infty}\left[x^3\left(-3+\frac1x\right)\right]
          =\left(\lim_{x\to+\infty}x^3\right) \cdot \lim_{x\to+\infty}\left(-3+\frac1x\right)\\[1em]
          &=\underbrace{\lim_{x\to+\infty}x^3}_{+\infty}
            \cdot\left(
              \underbrace{\lim_{x\to+\infty}(-3)}_{-3}
              +\underbrace{\lim_{x\to+\infty}\frac1x}_{0}
            \right)=-\infty.\quad
            \text{\raisebox{-.25\height}{\CachedTikzFigure{../_shared/resources/tikz/smiley_happy.tex}}}
      \end{aligned}
    \]
  \end{ClozeArea}
\end{ExampleBox}

\begin{KeyPointBox}{Unbestimmte Formen}{real_functions_indeterminate_forms}
  Unbestimmte Formen \EnglishTerm{indeterminate forms} sind als Rechenergebnisse zu vermeiden:
  Aus ihnen allein lässt sich kein Grenzwert bestimmen.
  \[
    \frac00,\qquad\frac\infty\infty,\qquad\infty-\infty,
    \qquad0\cdot\infty,\qquad1^\infty,\qquad0^0,\qquad\infty^0.
  \]
\end{KeyPointBox}


